\chapter{Myopia}

\medskip
\noindent\textit{\textbf{Under review.} This chapter is an AI-assisted educational draft; it is not a substitute for professional medical advice.}

\section*{Definition and Overview}
short-sightedness in which distant objects appear blurred because the eye focuses images in front of the retina. Individual assessment is important because symptoms and treatment vary with the cause.

\section*{Clinical Features}
Blurred distance vision, squinting, headaches, eye strain, and difficulty seeing classroom boards or road signs.

\section*{When to Seek Medical Care}
Seek medical review for new, persistent, worsening, or function-limiting symptoms. Seek emergency help for severe breathing difficulty, collapse, sudden neurological change, severe chest or abdominal pain, or any rapidly worsening illness.

\section*{Diagnosis and Investigations}
Visual-acuity and refraction testing, with examination of the eye; dilated retinal examination is important in high or rapidly progressive myopia.

\section*{Management}
Glasses or contact lenses; in selected adults, refractive surgery; in children, outdoor time and clinician-guided myopia-control options such as specialised lenses or low-dose atropine.

\section*{Complications}
Possible complications depend on the cause and may include progression of the condition, effects on other organs, treatment adverse effects, and reduced quality of life. Early assessment can reduce preventable harm.

\section*{Prognosis and Outcomes}
Outcome depends on the underlying cause, severity, complications, and response to treatment. Discuss individual prognosis with the treating clinician.

\section*{Prevention and Self-Care}
Follow the treatment plan, keep relevant appointments, avoid known triggers, protect affected areas from injury, and seek help early if symptoms worsen. Do not start or stop prescription treatment without clinical advice.

\clearpage
